\section{Results}\label{results}

\subsection{Foundation-specific probe accuracy}\label{foundation-specific-probe-accuracy}

All six foundation-specific probes achieve near-perfect accuracy
across the full depth of OLMo-2 1B. Care/harm and fairness/cheating
reach 100\% at nearly every layer; sanctity/degradation reaches
100\% at 12 of 16 layers. Liberty/oppression, loyalty/betrayal, and
authority/subversion show minor fluctuations (minimum 87.5\%, at
layers 3 and 15 for loyalty and layers 6 and 13--15 for authority).
The onset threshold (0.6) is exceeded at layer 0 for all
foundations.

These results confirm that moral foundation content is linearly
separable from the earliest layer --- consistent with the ``immediate
onset'' finding of \citet{reblitzrichardson2026fragility} for the
pooled binary moral/neutral probe. The per-foundation decomposition
reveals no qualitative difference in \emph{detectability} across
foundations: all are equally easy to decode. The interesting
variation is not in accuracy but in the \emph{geometry} of the probe
directions that achieve this accuracy.

On OLMoE-1B-7B, all foundations similarly reach near-perfect
accuracy, with authority/subversion showing the lowest peak (93.8\%
at layer 0). The accuracy profiles are essentially indistinguishable
between architectures.

\subsection{Framework geometry: integration, not collapse}\label{framework-geometry-integration-not-collapse}

\begin{figure}[t]
\centering
\includegraphics[width=\linewidth]{fig1_cosine_heatmap.pdf}
\caption{Pairwise cosine similarity between foundation probe directions at layer 0 (peak separation), OLMo-2 1B. Individualizing foundations (care, fairness, liberty) show higher within-group similarity than between-group pairs.}
\label{fig:cosine_heatmap}
\end{figure}

The headline finding: foundation probe directions are \emph{separated},
not collapsed. At the peak separation layer (layer 0), the mean
pairwise cosine similarity between the six foundation directions is
\textbf{0.272} --- far below the collapse threshold (\(>0.95\)) and below
the intermediate zone (\(0.8\)--\(0.95\)). The model does not encode
moral content through a single ``moral salience'' direction; it
maintains distinct directions for distinct moral foundations.

However, the cosine similarities are uniformly \emph{positive} (range
0.18--0.37 at layer 0), indicating that the foundation directions
share a common component. This is the \emph{integration} signature from
our trichotomy: the directions are separated but non-orthogonal,
consistent with a shared moral-salience subspace from which
foundation-specific directions deviate.

\textbf{Effective dimensionality.} The six foundation directions span
5 effective dimensions (the number of PCs explaining \(\geq 90\%\)
of variance) at every layer. This is near the maximum possible for
6 directions, confirming that the directions are geometrically
distinct --- they do not collapse into a lower-dimensional subspace.

\begin{figure}[t]
\centering
\includegraphics[width=\linewidth]{fig2_layerwise_geometry.pdf}
\caption{Layer-wise geometric metrics for OLMo-2 1B. (a) Mean pairwise cosine similarity peaks at middle layers. (b) Effective dimensionality remains constant at 5 across all layers.}
\label{fig:layerwise}
\end{figure}

\subsection{MFT group structure in the dendrogram}\label{mft-group-structure-in-the-dendrogram}

\begin{figure}[t]
\centering
\includegraphics[width=0.7\linewidth]{fig3_dendrogram.pdf}
\caption{Hierarchical clustering (Ward's method) of foundation probe directions at layer 0, OLMo-2 1B. The first split perfectly recovers the MFT individualizing/binding distinction.}
\label{fig:dendrogram}
\end{figure}

Hierarchical clustering of the six foundation directions at layer 0
produces a dendrogram that perfectly recovers the MFT
individualizing/binding distinction: the first split separates
\{liberty, care, fairness\} from \{sanctity, loyalty, authority\}.
Within the individualizing cluster, care and fairness merge first;
within the binding cluster, loyalty and authority merge first.

This alignment with MFT's predicted group structure is notable:
the model's representation geometry mirrors a theoretical
distinction from moral psychology that was not an explicit training
objective. The probe directions were trained independently for each
foundation; the clustering structure emerges from the representations
themselves.

The permutation test for the individualizing/binding distinction
reaches significance at layer 0 (\(p = 0.012\)) and layer 3
(\(p = 0.035\)), though not at most other layers (median \(p = 0.23\)).
With only 6 items and 20 unique 3--3 partitions, statistical power
is limited. The significant result at layer 0 corroborates the
dendrogram structure; the layer-to-layer variation in \(p\)-values
reflects the sensitivity of the scalar within/between summary
statistic to small changes in pairwise cosines.

\subsection{Layer-wise geometric development}\label{layer-wise-geometric-development}

The geometric structure evolves across layers in a characteristic
pattern. Mean pairwise cosine similarity is lowest at layer 0
(0.272), rises to a peak of 0.366 at layer 10, then gradually
decreases through the remaining layers (0.306 at layer 15). This
inverted-U pattern indicates that middle layers develop \emph{more}
shared structure between foundation directions (partial collapse),
while early and late layers maintain greater separation.

Effective dimensionality remains constant at 5 across all layers,
meaning the rank of the direction set does not change even as the
pairwise angles shift. The geometric development is a rotation of
the direction set within a fixed-rank subspace, not a
dimensionality change.

\subsection{Dense vs.~MoE framework geometry}\label{dense-vs.-moe-framework-geometry}

Framework geometry is remarkably similar between architectures.
OLMoE-1B-7B shows mean pairwise cosine similarity ranging from
0.260 to 0.350, compared to OLMo-2 1B's 0.275 to 0.363. Effective
dimensionality is 5 at all layers for both models. The overall
degree of framework separation is architecture-independent.

The dendrogram structure diverges at specific layers. At layer 7,
OLMoE produces a \emph{perfect} MFT split: \{liberty, care, fairness\}
vs.~\{loyalty, authority, sanctity\}. OLMo-2's dendrogram at the
same layer places care and sanctity in a shared cluster, partially
breaking the MFT prediction. The MFT group structure appears at
\emph{some} layer in both architectures, but the specific layer varies.

This finding extends \citet{reblitzrichardson2026dilution}: output
dilution affects moral encoding \emph{scale} (77\(\times\) signal gap) but
not framework \emph{structure}. The geometric organization of moral
foundations is preserved across architectures.

\subsection{Direction stability under bootstrap}\label{direction-stability-under-bootstrap}

Bootstrap resampling (200 iterations) reveals a stability gradient
across layers. Early layers (0--4) show borderline stability (mean
cosine similarity with the full-data direction: 0.73--0.80, below
the 0.8 threshold). Middle and late layers (5--14) are stable
(mean cosine \(> 0.80\)). Care/harm is the most stable foundation at
every layer.

This gradient has two implications. First, the geometric analysis
at early layers --- including the peak separation layer (layer 0)
--- should be interpreted with caution, as the specific pairwise
cosine values may shift under resampling. Second, the stability
gradient itself is informative: probe directions become more
determined as representations become more specialized, paralleling
the lexical-to-compositional gradient from
\citet{reblitzrichardson2026fragility}.

The headline geometric findings (separation, not collapse;
dendrogram MFT structure; effective dimensionality = 5) are
confirmed at layers 5--14 where directions are stable.

\subsection{Differential fragility across frameworks}\label{differential-fragility-across-frameworks}

\begin{figure}[t]
\centering
\includegraphics[width=\linewidth]{exp7_mean_critical_bars.pdf}
\caption{Mean critical noise per foundation for OLMo-2 1B (left) and OLMoE-1B-7B (right). Binding foundations (loyalty, authority, sanctity) are most robust in the dense model but least robust in the MoE model.}
\label{fig:fragility}
\end{figure}

Per-foundation fragility reveals a striking cross-architectural
pattern.

\textbf{OLMo-2 1B (dense).} Binding foundations are the \emph{most robust}:
loyalty/betrayal has the highest mean critical noise (6.21),
followed by authority/subversion (5.53). Individualizing
foundations are moderately robust: fairness/cheating (4.43),
liberty/oppression (4.24), care/harm (4.07). Sanctity/degradation
is the most fragile (3.85). The universality hypothesis (care/harm
as most robust) is \emph{not} supported.

\textbf{OLMoE-1B-7B (MoE).} The pattern reverses. Individualizing
foundations are more robust: care/harm (2.10), fairness/cheating
(1.96), liberty/oppression (1.90). Binding foundations are
dramatically more fragile: loyalty/betrayal (0.86),
sanctity/degradation (0.91), authority/subversion (0.64).

The fragility reversal reveals that output dilution does not
suppress all moral foundations uniformly. Binding foundations ---
which \citet{graham2013mft} characterize as more culturally
variable and group-oriented --- lose proportionally more robustness
under the MoE output scale gap. In OLMo-2, binding foundations
show a mean critical noise of 5.24 vs.~4.25 for individualizing
(ratio 1.23\(\times\)). In OLMoE, the ratio inverts: individualizing
mean 1.98 vs.~binding mean 0.80 (ratio 2.48\(\times\) in the opposite
direction). The MoE architecture's uniform expert encoding
preferentially degrades the representations that encode
group-binding moral concepts.

\subsection{Geometric trajectory during training}\label{geometric-trajectory-during-training}

\begin{figure}[t]
\centering
\includegraphics[width=\linewidth]{fig6_geometric_trajectory.pdf}
\caption{Geometric trajectory during OLMo-2 1B pre-training. (a) Mean cosine similarity stabilizes by step 2000. (b) Effective dimensionality remains at 5 throughout. (c) Accuracy continues climbing after geometry plateaus.}
\label{fig:trajectory}
\end{figure}

Framework geometry develops rapidly and stabilizes early. Tracking
mean pairwise cosine similarity (averaged over stable layers 5--14)
across 20 OLMo-2 1B checkpoints from step 0 to step 35,000:

\begin{itemize}
\item
  \textbf{Step 0.} Mean cosine similarity is 0.007 --- the six
  foundation directions are effectively orthogonal, as expected for
  probes trained on random representations. Mean accuracy is 0.510
  (chance).
\item
  \textbf{Step 1000.} Cosine similarity jumps to 0.176. Accuracy reaches
  0.734. The model has already begun developing shared moral-salience
  structure.
\item
  \textbf{Step 2000.} Cosine similarity reaches 0.382 --- within 5\% of
  its final value. Accuracy is 0.873, still 10 points below its
  peak.
\item
  \textbf{Steps 3000--15,000.} Cosine similarity plateaus at
  \(0.38\)--\(0.40\). Accuracy continues climbing from 0.922 to 0.970.
\item
  \textbf{Steps 20,000--35,000.} Cosine similarity drifts slightly
  downward (0.382 to 0.367). Accuracy stabilizes at 0.975--0.979.
\end{itemize}

Effective dimensionality remains constant at 5 from step 0 onward.
This is expected: random unit vectors in \(\mathbb{R}^{2048}\) are
nearly orthogonal with high probability, so even at initialization
the six probe directions span 5 effective dimensions. The
informative signal is cosine similarity, which transitions from
\(\approx 0\) (random) to \(\approx 0.4\) (mature structure) during the
first 2000 training steps.

The temporal dissociation --- framework geometry stabilizing at step
2000 while accuracy continues improving through step 25,000
--- extends the two-phase pattern identified by
\citet{reblitzrichardson2026fragility}: accuracy saturates early, but
fragility continues resolving. Here we add a third metric:
inter-framework \emph{structure} also stabilizes before inter-framework
\emph{discriminability} finishes developing. The model discovers the
geometric layout of moral concepts early and then spends the
remainder of training strengthening the representations within that
fixed layout.
